\section{Summary of Findings}
We set out to turn CUAD's clause-finding benchmark into something a person signing a contract could use: a system that finds the clauses, says how risky they are, and explains them in plain English. On a strict contract-level 80/20 split, the finished system reaches 85.65\% accuracy and a micro-F1 of 0.779 for finding which clauses are present, and a token-F1 of 0.764 for extracting the clause text when given the right window. It runs live in the \emph{Clausify} app.

Those are scores for separate parts, and the most useful thing this project can say is that the scores of the parts are not the score of the system. Answering our five research questions (Section~\ref{sec:rqs}) gave three negative results, and they are our main contribution:

\begin{itemize}[itemsep=3pt]
  \item \textbf{The windowed transformer does not beat the lexical baseline} (RQ1), and it does not lose for any of the reasons we could have blamed. Not the training budget: the full 53{,}662-window run moves micro-F1 by $-0.0004$. Not the seed: three runs put the gap at 9.5 standard deviations. Not the aggregation rule: six alternatives, each at its own best threshold, all stay behind. The threshold accounts for part of it, an un-chosen 0.5 costs 0.056 micro-F1, and that qualification belongs on the finding, but even the best cell we could construct remains below the baseline.
  \item \textbf{The summarizer does not summarize} (RQ3). It emits a category template verbatim on every input, and a constant lookup table matches its score against clause-specific references. This is a proof-of-concept classification-to-template component and we describe it that way everywhere.
  \item \textbf{The pipeline loses most of what its parts find} (RQ4). 21.7\% end-to-end survival against $\approx$64\% predicted, localized to a window-selection failure and a train/inference mismatch: both fixable without a larger model.
\end{itemize}

RQ5 gave the result we least expected and think is the most important for anyone building a system like this. Our evaluation score and our stated goal never matched: the AND-ensemble wins on micro-F1 but misses 36.9\% of High-risk clauses, while the OR rule, which comes last on micro-F1, misses only 14.2\%. Half of the clauses behind the overall score are Low-risk: dates, party names, governing law. A system built to warn a signer before signing was being chosen by a score dominated by the clauses that matter least.

Five lessons seem likely to apply beyond this project.

\textbf{First, a strong simple baseline deserves to be taken seriously, and beating it must be shown, not assumed.} When it is not beaten, the other possible explanations should be tested one at a time, not accepted or brushed aside.

\textbf{Second, a lead should be compared with its own noise before it is reported.} We measured three sources: the bootstrap over contracts (half-width about 0.008), the model seed (SD 0.0015) and the train/test split (SD 0.0063). Two of the three are larger than the $+0.0042$ lead we had used to rank a table.

\textbf{Third, an automatic score measures its reference sentences, not the task.} ROUGE-L 0.775 became 0.116 when only the reference sentences changed.

\textbf{Fourth, test the whole system, not just the parts}, and test it against the harm it is meant to prevent, not against an average that treats a governing-law clause and an unlimited indemnity as equally important.

\textbf{Fifth, limits are cheap to measure and expensive to assume.} A one-minute retrieval check (64\%) ruled out a design before it cost any training time. The same habit, applied to sparse attention, corrected \emph{us}: a Longformer training step does fit in this laptop's memory at batch size 1, which is the opposite of what we had first claimed.

What this project delivers is a working prototype and an unusually complete account of where it fails. Two gaps are large enough that we would not present the system as more than a prototype until they are closed: \textbf{no one with legal training has reviewed the risk table}, which is the main thing this work adds to CUAD, and \textbf{no user has ever been watched using the app}, so its main claim, that a report sorted by risk helps a non-lawyer, is untested. Both are listed first in the next section, and neither needs more computing power.

\section{Contributions to the Field}
The main contributions of this work are:
\begin{enumerate}[itemsep=2pt]
  \item \textbf{A first risk table for the CUAD categories.} Each of the 41 categories gets a High\slash Medium\slash Low level and a written reason, judged from the weaker party's side (8 High, 22 Medium, 11 Low), together with clear rules for how levels were given (Section~\ref{sec:riskcriteria}). We present it as a \emph{starting point for experts to check}, not as a proven legal-risk model: the three of us wrote it, none of us has legal training, and no legal professional has reviewed it yet (Section~\ref{sec:taxvalid}). It gives risk to a \emph{category}, not to the wording of one clause, which is why we call it category-level risk sorting throughout.
  \item \textbf{A tested result on finding clauses in long documents.} A TF--IDF baseline that reads the whole document (micro-F1 0.775) turned out to be hard to beat. Our window-based DistilBERT loses to it (0.692), and we checked the two obvious excuses: training on all 53{,}662 window examples instead of 24{,}000 changed the score by only $-0.0004$, and across three random seeds the gap is 9.2 standard deviations wide. An AND-ensemble comes first on paper (0.779), but the bootstrap over contracts and the variation across splits both show that its lead cannot be told apart from zero, so we report the ensemble for the kind of mistakes it makes, not for its rank.
  \item \textbf{A measurement of how errors add up.} Testing the three steps one after another, instead of separately, shows that only 22.0\% of real clauses get through the whole system, against about 65\% predicted from the separate scores. It also finds the two causes, a wrong choice of window and a mismatch between training and real use, and both can be fixed without a larger model.
  \item \textbf{A demonstration that a small set of template sentences turns summarizing into classification.} Testing the same 150 clauses against our 41 category templates and against clause-specific references drops ROUGE-L from 0.775 to 0.116, and the fine-tuned model is shown to output a template word for word for 100\% of inputs.
  \item \textbf{A simple safeguard: measure the limit before training.} We show that question-guided retrieval of windows caps the possible recall at 64\%, and we used that one measurement to reject the whole retrieval approach before spending any training time on it.
  \item \textbf{A system that can be rebuilt on a laptop.} All three models train on a single Apple-Silicon laptop. The split, the trained weights and a separate test notebook are all saved, so every number in this project can be produced again from disk.
  \item \textbf{Clausify}, a prototype web app that runs the trained models live and shows the results grouped by risk. It demonstrates the system but has not been tested with users, and Section~\ref{sec:demolimits} records what it does not do well.
\end{enumerate}

\section{Threats to Validity}
\label{sec:threats}
The results in Chapter~5 depend on choices that could change the conclusions if they were wrong. We list them under the four usual headings, and in each case we say whether we \emph{measured} the threat or only \emph{note} it.

\subsection{Internal validity: could our own procedure have produced these numbers?}
\textbf{Settings were never tuned} (Section~\ref{sec:valprotocol}). This removes the usual threat, since no setting was chosen by trying versions on the test set, but it creates the opposite one: the transformer competes with standard default settings against a baseline that has almost no settings to get wrong. Section~\ref{sec:threshold} shows this is a real effect: the unchosen 0.5 threshold costs the transformer 0.061 micro-F1, so part of the main gap comes from the threshold, not from the model. \emph{We think this is the most serious threat, and it limits our main finding.}

\textbf{An early test run of the summarizer tracked its loss on test-split clauses} (Section~\ref{sec:valprotocol}). No model was chosen from that run, and the final model was trained again without any evaluation set, so no reported score depends on it. But it was the wrong practice, and the ``validation'' loss it printed should be ignored. Figure~\ref{fig:sumloss} now shows the October 2026 re-training, which used no evaluation set.

\textbf{Some design choices were made after looking at the training data.} This is allowed, but the line needs care. The window size and step follow from the median clause length measured on training contracts. The 64\% retrieval limit, however, was measured on the whole dataset, test contracts included. That measurement ruled out a design rather than choosing a setting, and it was made before any model existed, so we judge the risk of leakage to be negligible. It is not zero, though, and we would make the split first if we did the work again.

\textbf{No seed variation for the span model and the summarizer.} Only the presence model was re-run with new seeds (Section~\ref{sec:seeds}). The span model and the summarizer were each trained again with the same seed and gave the same scores to within 0.001, so their seed variance is still unmeasured.

\subsection{Construct validity: do our metrics measure what we claim?}
\textbf{ROUGE-L does not measure summary quality here, and we can show it.} Against our 41 templates it gives 0.775; against clause-specific references it gives 0.116, and a fixed lookup table scores almost the same, 0.115 (Section~\ref{sec:summ}). We report both numbers because the first one on its own would give a false picture. We still have no measure of \emph{faithfulness} or readability (Section~\ref{sec:lit_faith}), so even 0.116 does not show that the output keeps the legal meaning.

\textbf{Micro-F1 does not measure harm to the user.} Section~\ref{sec:riskeval} shows the two can point in opposite directions: the setup that wins on micro-F1 misses 36.9\% of High-risk clauses, while a setup that comes last misses 14.2\%. Half of the clauses behind the overall score are Low-risk: dates, party names, governing law. Any claim like ``this system is 77.9\% good'' has this problem.

\textbf{Category-level risk is not clause-level risk} (Section~\ref{sec:risk}). Our levels say that a \emph{category} usually deserves attention; they say nothing about whether a particular cap is generous or harsh. Calling the output a risk assessment would claim too much, which is why we call it sorting by risk.

\textbf{The risk levels are our own.} The risk-weighted results use our unchecked choice of which categories are High, Medium and Low (Section~\ref{sec:taxvalid}); if a lawyer moves categories between levels, those numbers change too.

\textbf{Token-F1 against one marked span assumes the clause has clear edges.} Section~\ref{sec:span} notes that the edges of legal clauses are often debatable, and Section~\ref{sec:errorexamples} shows at least one ``false negative'' that looks like a disagreement with the labels rather than a model mistake.

\subsection{External validity, where do these results transfer?}
Every number comes from CUAD: \textbf{510 US business contracts taken from SEC filings, in English}. We have not tested a single contract from outside that dataset. So we have no evidence about: other legal traditions or countries, including \textbf{Bangladeshi contracts}, which are the setting that motivated this work; consumer or employment contracts, which are built differently and have a different weaker party; contracts not in English; contracts much shorter or much longer than CUAD's; and recent machine-written contracts. The risk table has one more threat of this kind: whether a non-compete can be enforced depends on the country or state, so a level that is right in one place may be wrong in another. Because CUAD comes from SEC filings, it also leans towards large-company contracts where both sides have lawyers, which is close to the opposite of the small businesses and individuals the system is meant for.

\subsection{Statistical conclusion validity: are the intervals trustworthy?}
\textbf{We measured three sources of variation}, and they differ in size: the bootstrap over contracts gives an interval half-width of about 0.008 (Section~\ref{sec:bootstrap}); the model-seed SD of the ensemble is 0.0015 over three runs (Section~\ref{sec:seeds}); and the split SD is 0.0063 over twenty splits (Section~\ref{sec:splitvar}). The bootstrap and the split variation are both larger than the ensemble's $+0.0042$ lead, which is why we call it a tie. The split variation is the largest, and we could measure it only for the baseline.

\textbf{Three seeds is a small sample} for estimating a standard deviation; the uncertainty around 0.0015 is itself large.

\textbf{Scores for single rare categories are not meaningful.} Five categories have at most seven test examples. Their F1 values are reported for completeness, not as estimates, and no claim in this project depends on any one of them.

\textbf{One split, one dataset.} The bootstrap resamples contracts within our test set; it cannot show how results would change on another dataset, and no interval in this project should be read as covering that.

\section{Limitations}
\label{sec:limitations}
Section~\ref{sec:threats} looked at the threats to each kind of validity. This section lists the limits of the system itself: what it does not do, as opposed to what could weaken the measurements. Each limit is described where it comes up in the text; these are pointers.

\textbf{No usability study} (Section~\ref{sec:usability}). Nobody outside the three of us has used Clausify. The project argues that a report sorted by risk is more useful to a non-lawyer than a plain list of clauses, but offers no evidence for it: no measure of task completion, understanding, trust or over-reliance. Since Section~\ref{sec:threshold} shows the confidence bar is too confident and Section~\ref{sec:summ} shows the plain-English sentence is the same for every clause of a category, over-reliance is exactly what a usability study would need to look for.

\textbf{Model size.} We used smaller models (DistilBERT, FLAN-T5-small) to fit a laptop's computing budget, so our scores are lower than the same method would reach with larger models on CUDA hardware.

\textbf{Training with one seed and one data budget, partly checked.} Section~\ref{sec:budget} reports what happened when the presence transformer was trained on all 53{,}662 window examples instead of a sample of 24{,}000, and Section~\ref{sec:seeds} reports the spread across three training seeds. The other two models, the span model and the summarizer, have each been trained twice with the same seed (matching to within 0.001) but never with a new seed, and on a sample of the data (8{,}000 of 11{,}156 and 4{,}000 of 11{,}156). We have not measured their seed variance at all, and their scores should be read with that in mind.

\textbf{The rarest categories are close to impossible to learn with this much data}, and Section~\ref{sec:raretail} shows the AND-ensemble makes this worse, not better: on the ten rarest categories the AND rule gains no F1 over the transformer alone, buying precision only with recall, and it scores zero on three categories that the transformer alone was partly finding. Since one of the rare categories is High-risk in our table, this limits the real system, not just the ranking.

\textbf{What the span test covers.} Section~\ref{sec:span} measures extraction inside windows that contain the answer, not the heavier full-document precision-at-recall test. Section~\ref{sec:e2e} closes part of that gap by measuring what gets through the real system, but it is not the CUAD paper's method, and we do not present the two as comparable.

\textbf{The summarizer does not summarize.} Section~\ref{sec:summ} shows this with numbers, not just as a warning: the fine-tuned model outputs one of 41 pre-written templates word for word for every input, and against clause-specific references it scores no better than the template alone with no model at all. It is a ``find the category, show its template'' component and should be described as one.

\textbf{The risk table has not been checked by anyone outside the team.} It was written by three computer science students, and no legal professional has reviewed it yet (Section~\ref{sec:taxvalid}). It is also set per category, not per clause: it flags how risky a category usually is, not how harsh the exact wording in front of the user is.

\textbf{One dataset, one country.} Every number in this project comes from CUAD: 510 US business contracts taken from SEC filings. We have not tested the system on a single contract from outside that dataset, so we have no evidence about how it handles other ways of writing contracts, and none at all about the Bangladeshi setting that motivated the work.

\section{Recommendations for Future Work}
\label{sec:futurework}
\begin{itemize}[itemsep=2pt]
  \item \textbf{Get the risk table reviewed.} The first thing we would do with another week: give \texttt{review/risk\_taxonomy\_review\_sheet.csv} to two or three law students or lawyers, run \texttt{review/agreement.py}, and report Cohen's or Fleiss' $\kappa$ and the categories they disagree on, whatever the result. It takes each person about half an hour, and it is the cheapest way for this project to gain trust.
  \item \textbf{Watch people use Clausify.} Five to ten people, each reviewing a contract with and without the app, measuring the time taken, the clauses correctly found and, most importantly, whether users trust the output too much. The project claims that a report sorted by risk helps a non-lawyer; nothing in it tests that claim yet.
  \item \textbf{Choose the setup by the goal that matters.} Define a risk-weighted goal (High-risk recall, or an F1 that counts High-risk clauses more), take a validation split from the 408 training contracts, choose the ensemble rule and each model's threshold on it, and report once on the test set. Section~\ref{sec:riskeval} suggests OR would win; choosing it from the evidence we already have would mean choosing on the test set.
  \item \textbf{Calibrate the confidence score.} Platt scaling or isotonic regression on a validation split would make the app's confidence bar mean what a user thinks it means (ECE is 0.205 now). No retraining is needed.
  \item \textbf{Let the span model say ``not here''.} Training with unanswerable examples, as in SQuAD~v2, would stop it from answering confidently from windows that cannot contain the clause, which is one of the two causes of the 22.0\% whole-system rate (Section~\ref{sec:e2e}).
  \item \textbf{Train the span model on the windows it will really see.} Placing the answer at different positions during training, instead of always 150 characters in, deals with the other cause.
  \item \textbf{Test summaries for faithfulness, not word overlap.} Use BERTScore, a check that the summary follows from the clause, a readability score, and a small human rating of whether the legal meaning is kept (Section~\ref{sec:lit_faith}).
  \item \textbf{Learned pooling over windows.} The test in Section~\ref{sec:pooling} tried only fixed rules for combining window scores; a learned rule that weighs windows by relevance (attention pooling) is the untested option most likely to help.
  \item \textbf{Cross-validation instead of one split.} The split is the largest uncertainty we measured (Section~\ref{sec:splitvar}), and we could only measure it for the baseline.
  \item \textbf{A summarizer that really summarizes.} Grow the 41 templates into a set of targets that differ from clause to clause, and fine-tune again. The 150-clause trial set in Section~\ref{sec:summ} is a working first version of both the labelling method and the test that would show success; scaling it to a few thousand clauses is the real path from picking templates to simplifying clauses.
  \item \textbf{Risk-aware ensemble rules.} Section~\ref{sec:raretail} suggests, but does not test, using the permissive OR rule for High-risk categories and the AND rule for the rest, so that recall is protected exactly where a missed clause costs the signer most. This needs its own held-out test: choosing the rule after seeing the test set would make the result worthless.
  \item \textbf{Meaning-based retrieval.} Our 64\% retrieval limit (Section~\ref{sec:ceiling}) was measured only with TF--IDF matching. A sentence-encoder retriever matches on meaning instead of shared words and might raise that limit enough to make a ``retrieve, then read'' design competitive; the measurement is cheap and we should have run it.
  \item \textbf{A trained sparse-attention baseline.} Section~\ref{sec:lit_long} argues against Longformer partly from arithmetic and partly from cost, and Section~\ref{sec:lfbench} measures the cost, but nobody has actually trained one on this task. On CUDA hardware, a Longformer baseline is the fair comparison that our argument currently replaces with reasoning.
  \item \textbf{Larger models.} Re-run the same system with DeBERTa-v3-base or large and FLAN-T5-base on CUDA hardware to close the gap to the published CUAD results.
  \item \textbf{Full-document span test.} Build the ranked whole-document test so that we can report precision at 80\% recall and compare directly with the CUAD paper.
  \item \textbf{Finding conflicts between clauses.} Bring the rule engine designed in an earlier phase of the project (for example, \emph{Exclusivity} against \emph{Most Favored Nation}) into the app.
  \item \textbf{Risk for each clause.} Move from risk per category to risk per clause by reading the extracted text itself: the amount of a cap, the length of a non-compete.
  \item \textbf{Contracts from outside CUAD.} Even three local English-language business contracts, run through the system and checked by hand, would tell us more about how it generalizes than any further tuning on CUAD, and it is the first step towards the Bangladeshi use this work was aimed at.
\end{itemize}

